%%%%%%%%%%%%%%%%%%%%%%%%%%%%%%%%%%%%%%%%%%%%%%%%%%%%%%%%%%%%%%%
%
% Final report document for F2 Practicals.
%
%%%%%%%%%%%%%%%%%%%%%%%%%%%%%%%%%%%%%%%%%%%%%%%%%%%%%%%%%%%%%%%
\documentclass[12pt, letterpaper]{article}
\usepackage[utf8]{inputenc}
\usepackage[T1]{fontenc}
\usepackage[margin=1in]{geometry}
\usepackage{microtype}
\usepackage{graphicx}
\usepackage{amsmath}
\usepackage{booktabs}
\usepackage[hidelinks]{hyperref}

\title{\textbf{Prueba%%%%%%%%%%%%%%%%%%%%%%%%%%%%%%%%%%%%%%%%%%%%%%%%%%%%%%%%%%%%%%%
%
% Final report document for F2 Practicals.
%
%%%%%%%%%%%%%%%%%%%%%%%%%%%%%%%%%%%%%%%%%%%%%%%%%%%%%%%%%%%%%%%
\documentclass[12pt, letterpaper]{article}
\usepackage[utf8]{inputenc}
\usepackage[T1]{fontenc}
\usepackage[margin=1in]{geometry}
\usepackage{microtype}
\usepackage{graphicx}
\usepackage{amsmath}
\usepackage{booktabs}
\usepackage[hidelinks]{hyperref}

\title{\textbf{F2 Practicals}}
\author{Nicolas Andres Ardila Gonzalez\thanks{Funded by the Overleaf team.}}
\date{July 2026}

\graphicspath{{images/}}

\setlength{\parindent}{0pt}
\setlength{\parskip}{0.75em}

\begin{document}

\maketitle
This document demonstrates basic LaTeX structure, text formatting, figures, lists, math, and tables.

\newpage

\begin{abstract}
This document presents a short introduction to the practical exercises and the structure used in the final report.

The abstract summarizes the key sections, including formatted text, figures, lists, mathematics, and tables.

A clean layout helps make the content easier to read and the technical examples simpler to follow.
\end{abstract}

\newpage
\tableofcontents

\newpage

\section{Customize text formats}

Some of the \textbf{greatest}
discoveries in \underline{science}
were made by \textbf{\textit{accident}}.

Some of the greatest \emph{discoveries} in science were made by accident.

\textit{Some of the greatest \emph{discoveries} in science were made by accident.}

\section{Add images or figures}

% The example figure is commented out because the image file is missing.
%\begin{figure}[h]
%    \centering
%    \includegraphics[width=0.75\textwidth]{Paper journal logo.png}
%    \caption{This is an example of figure}
%    \label{fig:placeholder}
%\end{figure}

\section{Using lists}

\begin{enumerate}
    \item This is the first entry in our list.
    \item The list numbers increase with each entry we add.
\end{enumerate}

\vspace{20pt}

\section{Math Mode}

In physics, the mass-energy equivalence is stated by the equation $E=mc^2$, discovered in 1895 by Albert Einstein.

This is also typeset in inline math mode as $E=mc^2$ and display math mode as:
\[
E = mc^2
\]

In natural units ($c = 1$), the formula expresses the identity:
\begin{equation}
E = m
\end{equation}

Subscripts in math mode are written as $a_b$ and superscripts as $a^b$. These can be combined and nested to write expressions such as:
\[
T^{i_1 i_2 \dots i_p}_{j_1 j_2 \dots j_q} = T(x^{i_1}, \dots, x^{i_p}, e_{j_1}, \dots, e_{j_q})
\]

We write integrals using $\int$ and fractions using $\frac{a}{b}$:
\[
\int_0^1 \frac{dx}{e^x} = \frac{e-1}{e}
\]

Lower case Greek letters are written as $\omega$, $\delta$, etc., while upper case Greek letters are written as $\Omega$, $\Delta$.

Mathematical operators are prefixed with a backslash, e.g. $\sin(\beta)$, $\cos(\alpha)$, $\log(x)$.

\subsection{First example}

The well-known Pythagorean theorem
\[x^2 + y^2 = z^2\]
was proved to be invalid for other exponents, meaning there are no integer solutions for $n > 2$ in:
\[x^n + y^n = z^n\]

\subsection{Second example}

This is a simple math expression $\sqrt{x^2+1}$ inside text.

\begin{displaymath}
\sqrt{x^2+1}
\end{displaymath}

\begin{equation*}
\sqrt{x^2+1}
\end{equation*}

\section{Basic document structure}

Paragraphs are separated by leaving a blank line between them, instead of using manual line breaks. This keeps the text readable and the layout consistent.

In normal text, use a blank line to start a new paragraph. Use line breaks such as \\\" only for short explicit breaks within the same paragraph when needed.

\section{Creating tables}

\begin{center}
\begin{tabular}{c c c}
 cell1 & cell2 & cell3 \\
 cell4 & cell5 & cell6 \\
 cell7 & cell8 & cell9
\end{tabular}
\end{center}

\begin{center}
\begin{tabular}{|c|c|c|}
 \hline
 cell1 & cell2 & cell3 \\
 cell4 & cell5 & cell6 \\
 cell7 & cell8 & cell9 \\
 \hline
\end{tabular}
\end{center}

% Create a table with cleaner formatting
\begin{center}
 \begin{tabular}{cccc}
 \toprule
 Col1 & Col2 & Col3 & Col4 \\
 \midrule
 1 & 6 & 87837 & 787 \\
 2 & 7 & 78 & 5415 \\
 3 & 545 & 778 & 7507 \\
 4 & 545 & 18744 & 7560 \\
 5 & 88 & 788 & 6344 \\
 \bottomrule
\end{tabular}
\end{center}

Table \ref{table:data} shows how to add a table caption and reference a table.
\begin{table}[h!]
\centering
\begin{tabular}{cccc}
 \toprule
 Col1 & Col2 & Col3 & Col4 \\
 \midrule
 1 & 6 & 87837 & 787 \\
 2 & 7 & 78 & 5415 \\
 3 & 545 & 778 & 7507 \\
 4 & 545 & 18744 & 7560 \\
 5 & 88 & 788 & 6344 \\
 \bottomrule
\end{tabular}
\caption{Table to test captions and labels.}
\label{table:data}
\end{table}

\end{document}
}}
\author{Nicolas Andres Ardila Gonzalez\thanks{Funded by the Overleaf team.}}
\date{July 2026}

\graphicspath{{images/}}

\setlength{\parindent}{0pt}
\setlength{\parskip}{0.75em}

\begin{document}

\maketitle
This document demonstrates basic LaTeX structure, text formatting, figures, lists, math, and tables.

\newpage

\begin{abstract}
This document presents a short introduction to the practical exercises and the structure used in the final report.

The abstract summarizes the key sections, including formatted text, figures, lists, mathematics, and tables.

A clean layout helps make the content easier to read and the technical examples simpler to follow.
\end{abstract}

\newpage
\tableofcontents

\newpage

\section{Customize text formats}

Some of the \textbf{greatest}
discoveries in \underline{science}
were made by \textbf{\textit{accident}}.

Some of the greatest \emph{discoveries} in science were made by accident.

\textit{Some of the greatest \emph{discoveries} in science were made by accident.}

\section{Add images or figures}

% The example figure is commented out because the image file is missing.
%\begin{figure}[h]
%    \centering
%    \includegraphics[width=0.75\textwidth]{Paper journal logo.png}
%    \caption{This is an example of figure}
%    \label{fig:placeholder}
%\end{figure}

\section{Using lists}

\begin{enumerate}
    \item This is the first entry in our list.
    \item The list numbers increase with each entry we add.
\end{enumerate}

\vspace{20pt}

\section{Math Mode}

In physics, the mass-energy equivalence is stated by the equation $E=mc^2$, discovered in 1895 by Albert Einstein.

This is also typeset in inline math mode as $E=mc^2$ and display math mode as:
\[
E = mc^2
\]

In natural units ($c = 1$), the formula expresses the identity:
\begin{equation}
E = m
\end{equation}

Subscripts in math mode are written as $a_b$ and superscripts as $a^b$. These can be combined and nested to write expressions such as:
\[
T^{i_1 i_2 \dots i_p}_{j_1 j_2 \dots j_q} = T(x^{i_1}, \dots, x^{i_p}, e_{j_1}, \dots, e_{j_q})
\]

We write integrals using $\int$ and fractions using $\frac{a}{b}$:
\[
\int_0^1 \frac{dx}{e^x} = \frac{e-1}{e}
\]

Lower case Greek letters are written as $\omega$, $\delta$, etc., while upper case Greek letters are written as $\Omega$, $\Delta$.

Mathematical operators are prefixed with a backslash, e.g. $\sin(\beta)$, $\cos(\alpha)$, $\log(x)$.

\subsection{First example}

The well-known Pythagorean theorem
\[x^2 + y^2 = z^2\]
was proved to be invalid for other exponents, meaning there are no integer solutions for $n > 2$ in:
\[x^n + y^n = z^n\]

\subsection{Second example}

This is a simple math expression $\sqrt{x^2+1}$ inside text.

\begin{displaymath}
\sqrt{x^2+1}
\end{displaymath}

\begin{equation*}
\sqrt{x^2+1}
\end{equation*}

\section{Basic document structure}

Paragraphs are separated by leaving a blank line between them, instead of using manual line breaks. This keeps the text readable and the layout consistent.

In normal text, use a blank line to start a new paragraph. Use line breaks such as \\\" only for short explicit breaks within the same paragraph when needed.

\section{Creating tables}

\begin{center}
\begin{tabular}{c c c}
 cell1 & cell2 & cell3 \\
 cell4 & cell5 & cell6 \\
 cell7 & cell8 & cell9
\end{tabular}
\end{center}

\begin{center}
\begin{tabular}{|c|c|c|}
 \hline
 cell1 & cell2 & cell3 \\
 cell4 & cell5 & cell6 \\
 cell7 & cell8 & cell9 \\
 \hline
\end{tabular}
\end{center}

% Create a table with cleaner formatting
\begin{center}
 \begin{tabular}{cccc}
 \toprule
 Col1 & Col2 & Col3 & Col4 \\
 \midrule
 1 & 6 & 87837 & 787 \\
 2 & 7 & 78 & 5415 \\
 3 & 545 & 778 & 7507 \\
 4 & 545 & 18744 & 7560 \\
 5 & 88 & 788 & 6344 \\
 \bottomrule
\end{tabular}
\end{center}

Table \ref{table:data} shows how to add a table caption and reference a table.
\begin{table}[h!]
\centering
\begin{tabular}{cccc}
 \toprule
 Col1 & Col2 & Col3 & Col4 \\
 \midrule
 1 & 6 & 87837 & 787 \\
 2 & 7 & 78 & 5415 \\
 3 & 545 & 778 & 7507 \\
 4 & 545 & 18744 & 7560 \\
 5 & 88 & 788 & 6344 \\
 \bottomrule
\end{tabular}
\caption{Table to test captions and labels.}
\label{table:data}
\end{table}

\end{document}
